\section{规划模型与执行模型解耦}

\subsection{为什么需要 Planner-Executor 分工}
讲者后半段强调了层级化思路：高层模型负责计划，低层模型负责执行。这样可以把复杂任务拆成局部可控子任务，降低一步步直接决策的噪声。

\begin{figure}[H]
\centering
\includegraphics[width=0.82\textwidth]{frame-06-planning.jpg}
\caption{视频约 00:27:40：讲者讨论让 planning model 先产出计划，再由执行模型落地动作}
\end{figure}

\begin{knowledgebox}{解耦后的收益}
\begin{itemize}
    \item 计划层可做全局目标约束；
    \item 执行层可针对当前页面做局部最优动作；
    \item 错误定位更清晰，便于调试和回滚。
\end{itemize}
\end{knowledgebox}

\subsection{层级接口如何设计}
一个可落地接口通常包含三部分：子目标描述、执行边界、完成判定。没有清晰接口时，规划层与执行层会相互覆盖，造成 ``计划写得好但执行跑偏'' 的现象。

\begin{importantbox}{层级接口最小字段}
\begin{enumerate}
    \item \textbf{subgoal}：当前阶段目标；
    \item \textbf{allowed actions}：允许动作集合；
    \item \textbf{stop condition}：完成或失败的判定标准；
    \item \textbf{fallback rule}：异常时返回上层的条件。
\end{enumerate}
\end{importantbox}

\subsection{何时该重新规划}
执行中一旦触发高置信度异常（元素不存在、页面跳转异常、预算超限），系统应及时回到规划层，而不是盲目继续下一个动作。

\begin{warningbox}{规划层常见失败}
\begin{itemize}
    \item 计划粒度过粗，执行层无法落地；
    \item 计划粒度过细，失去层级价值；
    \item 缺少 ``何时重规划'' 规则，导致错误持续扩散。
\end{itemize}
\end{warningbox}

\subsection{本章小结}
Planner-Executor 解耦是长链路 Agent 的关键工程模式。真正决定效果的是接口设计与重规划策略，而不仅是换一个更大模型。

